\documentclass[10pt,a4paper]{amsart}
\usepackage{amsmath,amssymb,amsthm,graphicx}
\usepackage{verbatim}
\newenvironment{reference}{\comment}{\endcomment}
%\newenvironment{reference}{}{}
\newenvironment{slogan}{\comment}{\endcomment}
\newenvironment{history}{\comment}{\endcomment}

% For commutative diagrams we use Xy-pic
\usepackage[all]{xy}

% We use 2cell for 2-commutative diagrams.
\xyoption{2cell}
% Theorem environments.
%
\theoremstyle{plain}
\newtheorem{theorem}[subsection]{Theorem}
\newtheorem{proposition}[subsection]{Proposition}
\newtheorem{lemma}[subsection]{Lemma}

\theoremstyle{definition}
\newtheorem{definition}[subsection]{Definition}
\newtheorem{example}[subsection]{Example}
\newtheorem{exercise}[subsection]{Exercise}
\newtheorem{situation}[subsection]{Situation}

\theoremstyle{remark}
\newtheorem{remark}[subsection]{Remark}
\newtheorem{remarks}[subsection]{Remarks}

\numberwithin{equation}{subsection}

% Macros
%
\def\lim{\mathop{\mathrm{lim}}\nolimits}
\def\colim{\mathop{\mathrm{colim}}\nolimits}
\def\Spec{\mathop{\mathrm{Spec}}}
\def\Hom{\mathop{\mathrm{Hom}}\nolimits}
\def\Ext{\mathop{\mathrm{Ext}}\nolimits}
\def\SheafHom{\mathop{\mathcal{H}\!\mathit{om}}\nolimits}
\def\SheafExt{\mathop{\mathcal{E}\!\mathit{xt}}\nolimits}
\def\Sch{\mathit{Sch}}
\def\Mor{\mathop{\mathrm{Mor}}\nolimits}
\def\Ob{\mathop{\mathrm{Ob}}\nolimits}
\def\Sh{\mathop{\mathit{Sh}}\nolimits}
\def\NL{\mathop{N\!L}\nolimits}
\def\CH{\mathop{\mathrm{CH}}\nolimits}
\def\proetale{{pro\text{-}\acute{e}tale}}
\def\etale{{\acute{e}tale}}
\def\QCoh{\mathit{QCoh}}
\def\Ker{\mathop{\mathrm{Ker}}}
\def\Im{\mathop{\mathrm{Im}}}
\def\Coker{\mathop{\mathrm{Coker}}}
\def\Coim{\mathop{\mathrm{Coim}}}

% Boxtimes
%
\DeclareMathSymbol{\boxtimes}{\mathbin}{AMSa}{"02}

%
% Macros for moduli stacks/spaces
%
\def\QCohstack{\mathcal{QC}\!\mathit{oh}}
\def\Cohstack{\mathcal{C}\!\mathit{oh}}
\def\Spacesstack{\mathcal{S}\!\mathit{paces}}
\def\Quotfunctor{\mathrm{Quot}}
\def\Hilbfunctor{\mathrm{Hilb}}
\def\Curvesstack{\mathcal{C}\!\mathit{urves}}
\def\Polarizedstack{\mathcal{P}\!\mathit{olarized}}
\def\Complexesstack{\mathcal{C}\!\mathit{omplexes}}
% \Pic is the operator that assigns to X its picard group, usage \Pic(X)
% \Picardstack_{X/B} denotes the Picard stack of X over B
% \Picardfunctor_{X/B} denotes the Picard functor of X over B
\def\Pic{\mathop{\mathrm{Pic}}\nolimits}
\def\Picardstack{\mathcal{P}\!\mathit{ic}}
\def\Picardfunctor{\mathrm{Pic}}
\def\Deformationcategory{\mathcal{D}\!\mathit{ef}}
\usepackage{cleveref}
\crefname{originaltheorem}{original theorem}{original theorems}
\Crefname{originaltheorem}{Original theorem}{Original theorems}
\makeatletter
\newlabel{theorem-colimit}{{09YQ}{source}{}{sourcepointer-1}{}}
\newlabel{theorem-colimit@cref}{{[originaltheorem][0][]09YQ}{[1][0][]source}{}{}{}}
\makeatother
\crefname{originallemma}{original lemma}{original lemmas}
\Crefname{originallemma}{Original lemma}{Original lemmas}
\makeatletter
\newlabel{derived-lemma-two-ss-complex-functor}{{015J}{source}{}{sourcepointer-2}{}}
\newlabel{derived-lemma-two-ss-complex-functor@cref}{{[originallemma][0][]015J}{[1][0][]source}{}{}{}}
\makeatother
\title{Cohomology of a filtered colimit of bounded below complexes of etale sheaves}
\author{The Stacks Project Authors}
\date{stacks-project:etale-cohomology.tex:local-snapshot-20261011; GNU FDL 1.2 or later}
\begin{document}
\maketitle
\noindent\textit{Editorial note.} All delivered mathematical text is the author's own, copied from the pinned Stacks Project chapter file etale-cohomology.tex; no statement or proof step was rewritten, reordered or removed.\par\medskip
\noindent\textit{Editorial note.} The proof is the author's complete argument: both spectral sequences are displayed, the degreewise input is applied, and the finiteness that legitimizes the termwise comparison is stated. The closing sentence 'Some details omitted' concerns that bookkeeping and is preserved verbatim in the counted span, together with the author's note that an alternative argument using stupid truncations avoids spectral sequences.\par\medskip
\noindent\textit{Editorial note.} Every labelled reference occurring inside the counted spans is carried in assembly.external\_references as a non-load-bearing navigation pointer (same-chapter ones at their definition site, cross-chapter ones at the invocation inside the counted span), so no load-bearing mathematics is replaced by a stub. No bibliography citation occurs inside the counted spans.\par\medskip
\noindent\textit{Editorial note.} The checked chapter list for same-name collisions is the set of stacks chapters staged in this repository: chow, cohomology, divisors, etale-cohomology, flat, more-algebra, more-morphisms, morphisms, spaces-more-morphisms, varieties together with this card's chapter etale-cohomology.\par\medskip
\noindent\textit{Editorial note.} The preamble is the chapter's own native preamble excerpt already used for this source: the Xy-pic setup (native preamble lines 14-18) and the native theorem declarations with the mathematical macros (native preamble lines 154-217).\par\medskip
\noindent\textit{Allowed standard prerequisites (editorial).}\begin{itemize}
\item Cohomology commutes with filtered colimits for sheaves on a limit of schemes. Same chapter, Theorem theorem-colimit. Supplies the degreewise colimit statement for the terms of the complexes.
\item The spectral sequence of a complex under a functor. Derived Categories, Lemma derived-lemma-two-ss-complex-functor. Supplies the two spectral sequences compared in the proof.
\end{itemize}\medskip
\ifdefined\hypertarget\hypertarget{sourcepointer-1}{}\fi
\noindent\textit{Source reference.} Original theorem 09YQ: sources/stacks-etale-cohomology/pinned-original-source.tex, source lines 7400--7400. Navigation pointer invoked at proof line 7802; the counted span contains the pointer only and the cited result is not reproduced inside it, so the pointer is registered as a non-load-bearing navigation stub in the established form for Stacks material. No bibliography citation is involved.\par\medskip
\ifdefined\hypertarget\hypertarget{sourcepointer-2}{}\fi
\noindent\textit{Source reference.} Original lemma 015J: sources/stacks-etale-cohomology/pinned-original-source.tex, source lines 7819--7819. Navigation pointer invoked at proof line 7819; the counted span contains the pointer only and the cited result is not reproduced inside it, so the pointer is registered as a non-load-bearing navigation stub in the established form for Stacks material. No bibliography citation is involved.\par\medskip
\begin{lemma}
\label{lemma-colimit-variant-complexes}
Let $X = \lim_{i \in I} X_i$ be a limit of a directed system of schemes
with affine transition morphisms $f_{i'i} : X_{i'} \to X_i$. We assume
that $X_i$ is quasi-compact and quasi-separated for all $i \in I$.
Let $\mathcal{F}_i^\bullet$ be a complex of abelian sheaves on
$X_{i, \etale}$. Let $\varphi_{i'i} : f_{i'i}^{-1}\mathcal{F}_i^\bullet \to
\mathcal{F}_{i'}^\bullet$ be a map of complexes on $X_{i, \etale}$
such that $\varphi_{i''i} = \varphi_{i''i'} \circ f_{i'' i'}^{-1}\varphi_{i'i}$
whenever $i'' \geq i' \geq i$. Assume there is an integer $a$ such that
$\mathcal{F}_i^n = 0$ for $n < a$ and all $i \in I$.
Then we have
$$
H^p_\etale(X, \colim f_i^{-1}\mathcal{F}_i^\bullet) =
\colim H^p_\etale(X_i, \mathcal{F}^\bullet_i)
$$
where $f_i : X \to X_i$ is the projection.
\end{lemma}

\begin{proof}
This is a consequence of Theorem \ref{theorem-colimit}. Set
$\mathcal{F}^\bullet = \colim f_i^{-1}\mathcal{F}_i^\bullet$.
The theorem tells us that
$$
\colim_{i \in I} H_\etale^p(X_i, \mathcal{F}_i^n) =
H_\etale^p(X, \mathcal{F}^n)
$$
for all $n, p \in \mathbf{Z}$. Let us use the spectral sequences
$$
E_{1, i}^{s, t} = H_\etale^t(X_i, \mathcal{F}_i^s) \Rightarrow
H_\etale^{s + t}(X_i, \mathcal{F}_i^\bullet)
$$
and
$$
E_1^{s, t} = H_\etale^t(X, \mathcal{F}^s) \Rightarrow
H_\etale^{s + t}(X, \mathcal{F}^\bullet)
$$
of Derived Categories, Lemma \ref{derived-lemma-two-ss-complex-functor}.
Since $\mathcal{F}_i^n = 0$ for $n < a$ (with $a$ independent of $i$)
we see that only a fixed finite number of terms $E_{1, i}^{s, t}$
(independent of $i$) and $E_1^{s, t}$ contribute to
$H^q_\etale(X_i, \mathcal{F}_i^\bullet)$ and
$H^q_\etale(X, \mathcal{F}^\bullet)$ and $E_1^{s, t} = \colim E_{i, i}^{s, t}$.
This implies what we want. Some details omitted.
(There is an alternative argument using ``stupid'' truncations
of complexes which avoids using spectral sequences.)
\end{proof}

\end{document}
